\documentclass[12pt, a4paper]{article}
\usepackage[utf8]{inputenc}
\usepackage[albanian]{babel}
\usepackage{amsmath}
\usepackage{amsfonts}
\usepackage{graphicx}
\usepackage{booktabs}
\usepackage{geometry}
\geometry{margin=1in}

\title{\textbf{Raport Shkencor: Analiza e Modelit Epidemiologjik SEIRV me Reagim Kolektiv}}
\author{\textbf{Azem Elezi}}
\date{\today}

\begin{document}

\maketitle

\section{Hyrje dhe Formulimi Matematik}
Ky projekt studion dinamikën e përhapjes së një infeksioni duke përdorur modelin e zgjeruar \textbf{SEIRV} (Susceptible, Exposed, Infectious, Recovered, Vaccinated). Inovacioni kryesor i këtij modeli qëndron në përfshirjen e \textbf{sjelljes kolektive ($\eta$)}, ku shkalla e kontaktit bazë ($\beta_0$) nuk është konstante, por ndryshon në mënyrë dinamike si reagim ndaj numrit të të infektuarve aktivë ($I$) në popullatë.

Ekuacioni i transmetimit efektiv përcaktohet si:
\begin{equation}
\beta_{\text{eff}}(I) = \max\left(\beta_0 \left(1 - \eta \frac{I}{N}\right), 0\right)
\end{equation}

Sistemi i ekuacioneve diferenciale jolineare që rregullon modelin është:
\begin{align}
\frac{dS}{dt} &= -\beta_{\text{eff}}(I) \frac{S I}{N} - \nu S \\
\frac{dE}{dt} &= \beta_{\text{eff}}(I) \frac{S I}{N} - \sigma E \\
\frac{dI}{dt} &= \sigma E - \gamma I \\
\frac{dR}{dt} &= \gamma I \\
\frac{dV}{dt} &= \nu S
\end{align}

Ku parametrat fillestarë të sistemit janë përcaktuar si:
\begin{itemize}
    \item $N$: Popullata totale ($1,000,000$)
    \item $\beta_0$: Shkalla bazë e transmetimit ($0.5$)
    \item $\sigma$: Shkalla e kalimit nga i ekspozuar në të infektuar ($1/5 \text{ ditë}^{-1}$)
    \item $\gamma$: Shkalla e shërimit ($1/10 \text{ ditë}^{-1}$)
    \item $\nu$: Shkalla e vaksinimit ditor ($0.001$)
\end{itemize}

\section{Rezultatet dhe Analiza e Sensitivitetit}
Për të vlerësuar ndikimin e ndryshimit të sjelljes njerëzore (distancimi social, ndërgjegjësimi higjienik), u krye një analizë sensitiviteti duke skanuar parametrin $\eta$ në katër skenarë të ndryshëm: $\eta = [0.0, 0.4, 0.8, 1.2]$. Rezultatet numerike të simuluara nëpërmjet paketës \texttt{SciPy} janë përmbledhur në Tabelën \ref{tab:metrikat}.

\begin{table}[htbp]
\centering
\caption{Metrikat e Kulmit Epidemik (Peak Analysis) për vlera të ndryshme të $\eta$.}
\label: \label{tab:metrikat}
\vspace{0.5em}
\begin{tabular}{cccc}
\toprule
\textbf{Vlera e $\eta$} & \textbf{Kulmi i të Infektuarve ($I_{\max}$)} & \textbf{Dita e Kulmit (Peak Day)} & \textbf{Statusi i Kurbës} \\
\midrule
\textbf{0.0} (Pa reagim)     & 275,542 & 76.4 & Kurba më e lartë \\
\textbf{0.4} (Reagim i dobët)  & 263,450 & 76.9 & Fillimi i sheshimit \\
\textbf{0.8} (Reagim mesatar) & 250,561 & 77.4 & Zhvendosje në kohë \\
\textbf{1.2} (Reagim i fortë)  & 236,730 & 78.0 & Kurba më e sheshtë \\
\bottomrule
\end{tabular}
\end{table}

\section{Diskutimi i Grafikëve}
Nga simulimet e kryera (të ruajtura automatikisht në folderin \texttt{results/figures/}), vërehen dy dukuri kryesore epidemiologjike:
\begin{enumerate}
    \item \textbf{Efekti ``Sheshimi i Kurbës'' (Flattening the Curve):} Me rritjen e parametrit të sjelljes kolektive nga $\eta = 0.0$ në $\eta = 1.2$, numri maksimal i të infektuarve aktivë ulet nga $275,542$ në $236,730$ (një rënie prej rreth $14\%$). Kjo vërteton matematikisht se ndryshimi i sjelljes redukton ngarkesën e menjëhershme mbi sistemin spitalor.
    \item \textbf{Shtyrja e Kohës së Kulmit (Peak Delay):} Rritja e reagimit kolektiv e zhvendos ditën e pikut nga dita $76.4$ në ditën $78.0$. Kjo shtyrje kohore u jep autoriteteve shëndetësore dritare kritike për alokimin e burimeve mjekësore dhe avancimin e fushatës së vaksinimit ($V$).
\end{enumerate}

\section{Konkluzione}
Modeli i zhvilluar dëshmon me sukses se ecuria e një epidemie nuk varet vetëm nga vetitë biologjike të patogjenit ($\beta_0, \gamma$), por në mënyrë kritike nga \textbf{reagimi socio-sjellës i popullatës}. Ndërsa vaksinimi ($V$) siguron stabilizimin e sistemit në afatgjatë, sjellja kolektive dinamike ($\eta$) shërben si mekanizmi kryesor për menaxhimin dhe zbutjen e fazës akute të krizës.

\end{document}
